\documentclass[10pt,a4paper]{amsart}
\usepackage[margin=19mm]{geometry}
\usepackage{amsmath,amssymb,mathtools}
\usepackage[scr=rsfs,cal=euler]{mathalfa}
\usepackage{xfrac}
\usepackage[unicode,hidelinks,bookmarks=false]{hyperref}
\newcommand{\ie}{i.e.\ }
\newcommand{\cf}{cf.\ }
\newcommand{\resp}{respectively\ }
\newcommand{\Subsection}{\subsection}
\providecommand{\nobreakdash}{\nobreak\hbox{-}}
\setlength{\emergencystretch}{2em}
\multlinegap0pt
\usepackage{bm}
\newtheorem{defn}{Definition}
\newtheorem{thm}{Theorem}
\newtheorem{asm}{Assumption}
\newtheorem{lem}{Lemma}
\usepackage{cleveref}
\crefname{defn}{Definition}{Definitions}
\crefname{thm}{Theorem}{Theorems}
\crefname{asm}{Assumption}{Assumptions}
\crefname{lem}{Lemma}{Lemmas}
\newcommand{\Brac}[1]{\left(#1\right)}
\newcommand{\E}{\mathbb{E}}
\newcommand{\Fro}[1]{\norm{#1}_{\rm F}}
\newcommand{\Rectbrac}[1]{\left[#1\right]}
\newcommand{\bX}{\boldsymbol{X}}
\newcommand{\bx}{{\boldsymbol{x}}}
\newcommand{\gto}{\rightarrow}
\newcommand{\mN}{\mathbb{N}}
\newcommand{\mR}{\mathbb{R}}
\newcommand{\mZ}{\mathbb{Z}}
\newcommand{\matT}{^{\!\top}}
\newcommand{\mcN}{\mathcal{N}}
\newcommand{\mcP}{\mathcal{P}}
\newcommand{\mcX}{\mathcal{X}}
\newcommand{\mdd}{\mathrm{d}}
\newcommand{\norm}[1]{\left\|#1\right\|}
\newcommand{\normo}[1]{\|#1\|}
\newcommand{\sfC}{\mathsf{C}}
\newcommand{\sfG}{\mathsf{G}}
\newcommand{\sfI}{\mathsf{I}}
\newcommand{\sfK}{\mathsf{K}}
\newcommand{\sfL}{\mathsf{L}}
\newcommand{\sfR}{\mathsf{R}}
\newcommand{\sfU}{\mathsf{U}}
\newcommand{\sfV}{\mathsf{V}}
\newcommand{\sfW}{\mathsf{W}}
\newcommand{\tr}{\text{\rm tr}}
\title{Stein's Method for Marginals on Large Graphical Models}
\author{Tiangang Cui, Shuigen Liu, Xin T. Tong}
\date{Source: arXiv:2410.11771v4 (CC BY 4.0)}
\begin{document}
\maketitle

\noindent\emph{Source and licence.} Stein's Method for Marginals on Large Graphical Models. Version arXiv:2410.11771v4 (CC BY 4.0), distributed by arXiv under the Creative Commons Attribution 4.0 International licence, http://creativecommons.org/licenses/by/4.0/, as recorded in the arXiv licence metadata for this exact version and shown on the abstract page https://arxiv.org/abs/2410.11771v4. Full licence notice: \texttt{licenses/arxiv-2410.11771-CC-BY-4.0.txt}.\par\smallskip

\noindent\textit{Editorial scope.} Counted unit: the article's thm thm:margin (source0.tex lines 601--612). The article's complete original proof of it is delivered with it (source0.tex lines 2025--2153, one \texttt{proof} environment of 129 lines, which the article places away from the statement). No mathematics is written, repaired or re-derived: the statement and the proof are the article's own lines, in the article's own notation, and no retained line is substituted or re-flowed.  Retained uncounted as the source-local context the proof needs: lines 186--200; lines 378--386; lines 420--427; lines 744--760; lines 1974--1981.  Nothing was omitted from any retained span, so the package carries no omission record.  No retained line contains a citation, so the wrapper carries no bibliography and no bibliography entry.  No retained line contains a figure environment, an image-inclusion command or drawing code, so no image is delivered and the package declares no asset dependency.  Editorial formatting only, in addition: an AMS \texttt{amsart} preamble that restates the article's own declarations and macros used by the retained lines, namely the environments \texttt{defn}, \texttt{thm}, \texttt{asm}, \texttt{lem} on separate counters, together with the standard packages those lines need.  The retained environments print in this document as Definition 1, Theorem 1, Asm 1, Lemma 1 in order of appearance, whereas the article prints the same items with the numbering of its own full document.  The complete native source of the article as distributed in the arXiv e-print for this exact version is bundled unchanged as \texttt{sources/166883/source0.tex}.
\begin{defn}[$\delta$-locality]
\label{def:PoisGrad}
Let $\pi \in \mcP(\mR^d)$ have a positive $C^1$ density and satisfy a Poincar\'e inequality.
We say that $\pi$ is {\bf $\delta$-localized} for some dimension-free constant $\delta>0$ if, for every $i\in[K]$ and every $1$-Lipschitz function $\phi : \mR^{d_i} \gto \mR$, the solution $u(\bx)$ to the \emph{marginal Stein equation}
\begin{equation}
\label{eqn:PoisEqn}
    \sfL_\pi u (\bx) := \Delta u(\bx) + \nabla \log \pi(\bx) \cdot \nabla u(\bx) = \phi(\bx_i) - \E_\pi [\phi(\bx_i)]
\end{equation}
exists in $H_0^1(\pi)$ and satisfies the gradient bound
\begin{equation}
\label{eqn:PoisGradBound}
    \norm{ \nabla u }_{\infty,1} := \sum_{j=1}^K \norm{ \nabla_j u }_{L^\infty} \le \delta.
\end{equation}
Here $\norm{ \nabla_j u }_{L^\infty} := \sup_{\bx\in\mR^d} \norm{ \nabla_j u(\bx) }$ denotes the $L^\infty$ norm of $\norm{\nabla_j u}_2$.
\end{defn}

\begin{defn}
\label{def:loc_graph}
An undirected graph $\sfG$ is called $\boldsymbol{(S,\nu)}$\textbf{-local} if there exist $S,\nu \in \mZ_+$ such that the cardinality of the order-$q$ neighbor set satisfies
\begin{equation}
\label{eqn:loc_graph}
    |\mcN(j;q)| \le 1 + S q^\nu , \quad \forall j \in \sfV, ~ \forall q\in\mN,
\end{equation}
where $S$ controls the maximal size of immediate neighborhoods, and $\nu$ is the effective ambient dimension of the graph that controls the growth of the neighborhood volume as a function of the radius $q$.
\end{defn}

\begin{thm}
\label{thm:MRF_loc}
Suppose $\pi \in \mcP(\mR^d)$ is associated with a $(S,\nu)$-local dependency graph $\sfG$. Suppose further there exist constants $m,M$ with $0<m\le M<\infty$ such that the local precision matrix of the density $\pi$ satisfies
\(
    m I_d \preceq - \nabla^2 \log \pi(\bx) \preceq M I_d, \forall \bx \in \mR^d.
\)
Then, the density $\pi$ is $\delta$-localized with $\delta = \frac{S \nu! \kappa^\nu }{m}$ where $\kappa = \frac{M}{m}$ is the condition number. 
\end{thm}

\begin{asm}
\label{asm:LLIS}
To facilitate the construction of local likelihood-informed subspaces, we partition the parameter space $\mcX = \mR^d$ into $K$ disjoint blocks, $\bx = (\bx_j)_{j\in[K]},$ where $\bx_j \in \mR^{d_j},\, \sum_{j=1}^K d_j = d$,
and assume the target distribution $\pi$ satisfies the following conditions.
\begin{itemize}
    \item The prior $\pi_0$ and the posterior $\pi$ share the same dependency graph $\sfG$; that is, for any $i,j\in\sfV$ that are not adjacent in $\sfG$, their mixed partial derivatives vanish:
    \begin{equation}
        \nabla_{ij}^2 \log \pi_0 (\bx) = \nabla_{ij}^2 \log \pi (\bx) = 0, \quad \forall i \notin \mcN(j),~ \forall j \in \sfV.
    \end{equation}
    Moreover, $\sfG$ is $(S,\nu)$-local (see \Cref{def:loc_graph}).
    \item There exist constants $m,M > 0$ such that for all $\bx \in \mR^d$,
    \begin{equation}
        m I_d \preceq - \nabla^2 \log \pi_0(\bx) \preceq M I_d, \qquad m I_d \preceq - \nabla^2 \log \pi(\bx) \preceq M I_d.
    \end{equation}
    Here $m,M$ are the same constants as those in \Cref{thm:MRF_loc} which guarantee that the local precision matrices are well conditioned.
\end{itemize}
\end{asm}

\begin{lem}
\label{lem:Kl1Bound}
Under \Cref{asm:LLIS}, and denote $\sfK = \sfC_{r,r}^\dag \sfC_{r,\bot}$, where $\sfC_{\bot,r} = \Pi_{\bot} \sfC \Pi_{r}$ and $\sfC_{r,r} = \Pi_{r} \sfC \Pi_{r}$
Then it holds that 
\[
    \sum_{k\in[K]} \norm{ \sfK(j,k) } \le S(1+S) \nu! \kappa^{\nu+2}. 
\]
\end{lem}

\begin{thm}
\label{thm:margin}
Consider a target $d$-dimensional random vector $\bX$ with a density $\pi \in \mcP_1(\mR^d)$, partitioned into $K$ disjoint subsets, $\bX = (\bX_j)_{j \in [K]}$.
%
Suppose $\pi'\in \mcP_1(\mR^d)$ is a $\delta$-localized approximation to $\pi$ in the sense of \Cref{def:PoisGrad}.
Then, the maximal $\sfW_1$ distance between marginal distributions of $\pi,\pi'$ satisfies
\begin{equation}
\label{eqn:MarginIneqn}
    \max_{i\in [K]} \sfW_1 (\pi_i^{}, \pi_i') \le \delta \cdot \max_{j\in [K]} \norm{ \nabla_j \log \pi' - \nabla_j \log \pi }_{L^1(\pi)},
\end{equation}
where $\pi_i$ and $\pi_i'$ denote the marginal densities of $\pi$ and $\pi'$ on the block of variables $\bX_i$, respectively. 
\end{thm}

\begin{proof}[Proof of Theorem~6]

By \Cref{asm:LLIS} and \Cref{thm:MRF_loc}, $\pi$ is $\frac{S\nu! \kappa^\nu}{m}$-localized.
Therefore, by \Cref{thm:margin}, it holds that 
\[
    \max_{i \in [K]} \sfW_1(\pi_i,\pi_i') \le \frac{S\nu! \kappa^\nu}{m} \cdot \max_{j \in [K]} \norm{ \nabla_j \log \pi - \nabla_j \log \pi' }_{L^1(\pi')}.
\]
Recall that $\pi(\bx) \propto \ell(\bx) \pi_0(\bx)$ and $\pi'(\bx) \propto \ell_r(\bx_r) \pi_0(\bx)$, where
\[
    \log \ell_r(\bx_r) = \int \log \ell(\bx_r,\widetilde{\bx}_\bot) \pi_0(\widetilde{\bx}_\bot|\bx_r) \mdd \widetilde{\bx}_\bot.
\]
Because $\Pi_{j,\bot}\nabla_j \log \ell_r(\bx_r)=0$, we have
\begin{align*}
    \nabla_j \log \pi (\bx) \,-\, &\nabla_j \log \pi' (\bx) = \nabla_j \log \ell (\bx) - \nabla_j \log \ell_r (\bx_r) \\
    =~& \Pi_{j,\bot} \nabla_j \log \ell (\bx) + \Pi_{j,r} \nabla_j \log \ell (\bx) - \int \Pi_{j,r} \nabla_j \log \ell(\bx_r,\widetilde{\bx}_\bot) \pi_0(\widetilde{\bx}_\bot|\bx_r) \mdd \widetilde{\bx}_\bot \\
    &- \int \log \ell(\bx_r,\widetilde{\bx}_\bot) \Pi_{j,r} \nabla_j \pi_0(\widetilde{\bx}_\bot|\bx_r) \mdd \widetilde{\bx}_\bot. 
\end{align*}
Thus
\begin{equation}
\label{eqn:pf_LLIS_err_decomp}
\begin{split}
    &\norm{ \nabla_j \log \pi - \nabla_j \log \pi' }_{L^1(\pi')} \\
    \le~& \int \norm{ \Pi_{j,\bot} \nabla_j \log \ell (\bx) } \pi'(\bx) \mdd \bx \\
    &+ \int \norm{ \Pi_{j,r} \nabla_j \log \ell (\bx) - \int \Pi_{j,r} \nabla_j \log \ell(\bx_r,\widetilde{\bx}_\bot) \pi_0(\widetilde{\bx}_\bot|\bx_r) \mdd \widetilde{\bx}_\bot } \pi'(\bx) \mdd \bx \\
    &+ \int \norm{ \int \log \ell(\bx_r,\widetilde{\bx}_\bot) \Pi_{j,r} \nabla_j \pi_0(\widetilde{\bx}_\bot|\bx_r) \mdd \widetilde{\bx}_\bot } \pi'(\bx) \mdd \bx \\
    =:& \sfI_1 + \sfI_2 + \sfI_3.
\end{split}
\end{equation}
Next we control the three terms separately.

\noindent
(1) For $\sfI_1$, by the Cauchy--Schwarz inequality,
\begin{equation}
\label{eqn:pf_LLIS_I1}
\begin{split}
    \sfI_1^2 \le~& \int \norm{ \Pi_{j,\bot} \nabla_j \log \ell (\bx) }^2 \pi'(\bx) \mdd \bx \\
    =~& \int \tr \Rectbrac{  \Pi_{j,\bot} \nabla_j \log \ell (\bx) \nabla_j \log \ell(\bx) \matT \Pi_{j,\bot}\matT } \pi'(\bx) \mdd \bx \\
    =~& \tr \Rectbrac{ \Pi_{j,\bot} \Brac{  \int \nabla_j \log \ell(\bx) \nabla_j \log \ell(\bx) \matT \pi'(\bx) \mdd \bx } \Pi_{j,\bot}\matT } \\
    =~& \tr \Rectbrac{ \Pi_{j,\bot} \sfR_j \Pi_{j,\bot}\matT }.
\end{split}
\end{equation}
Recall that $\sfR_j := \int \nabla_j \log \ell(\bx) \nabla_j \log \ell(\bx) \matT \pi'(\bx) \mdd \bx$.

\ \\[-5pt]
\noindent
(2) For $\sfI_2$, similarly we have
{\small \begin{align*}
    \sfI_2^2 \le~& \int \norm{ \Pi_{j,r} \nabla_j \log \ell (\bx) - \int \Pi_{j,r} \nabla_j \log \ell(\bx_r,\widetilde{\bx}_\bot) \pi_0(\widetilde{\bx}_\bot|\bx_r) \mdd \widetilde{\bx}_\bot }^2 \pi'(\bx) \mdd \bx  \\
    =~& \int \Brac{ \int  \norm{ \Pi_{j,r} \nabla_j \log \ell (\bx) - \int \Pi_{j,r} \nabla_j \log \ell(\bx_r,\widetilde{\bx}_\bot) \pi_0(\widetilde{\bx}_\bot|\bx_r) \mdd \widetilde{\bx}_\bot }^2 \pi_0(\bx_\bot|\bx_r) \mdd \bx_\bot } \pi'(\bx_r) \mdd \bx_r.  
\end{align*}}
As a conditional distribution of $\pi_0$, $\pi_0(\cdot|\bx_r)$ is $m$-log-concave. By the Poincar\'e inequality,
\begin{align*}
    \int \norm{ \Pi_{j,r} \nabla_j \log \ell (\bx) - \int \Pi_{j,r} \nabla_j \log \ell(\bx_r,\widetilde{\bx}_\bot) \pi_0(\widetilde{\bx}_\bot|\bx_r) \mdd \widetilde{\bx}_\bot }^2 \pi_0(\bx_\bot|\bx_r) \mdd \bx_\bot & \\
    \le \frac{1}{m} \int \Fro{ \nabla_{\bx_\bot} \Brac{ \Pi_{j,r} \nabla_j \log \ell (\bx)}\matT  }^2 \pi_0(\bx_\bot|\bx_r) \mdd \bx_\bot& \\
    = \frac{1}{m} \int \Fro{ \Pi_\bot \nabla \nabla_j\matT \log \ell (\bx) \; \Pi_{j,r}\matT }^2 \pi_0(\bx_\bot|\bx_r) \mdd \bx_\bot&.
\end{align*}
Substituting this into the above inequality, we have 
\begin{equation}
\label{eqn:pf_LLIS_I2}
\begin{split}
    \sfI_2^2 \le~& \frac{1}{m} \int \Brac{ \int \Fro{ \Pi_\bot \nabla \nabla_j\matT \log \ell (\bx) \; \Pi_{j,r}\matT }^2 \pi_0(\bx_\bot|\bx_r) \mdd \bx_\bot } \pi'(\bx_r) \mdd \bx_r \\
    =~& \frac{1}{m} \int \Fro{ \Pi_\bot \nabla \nabla_j\matT \log \ell (\bx) \; \Pi_{j,r}\matT }^2 \pi'(\bx) \mdd \bx \\
    =~& \frac{1}{m} \int \tr \Rectbrac{ \Pi_\bot \nabla \nabla_j\matT \log \ell(\bx) \, \Pi_{j,r}\matT \Pi_{j,r} \, \nabla_j \nabla\matT \log \ell(\bx) \Pi_\bot\matT } \pi'(\bx) \mdd \bx \\
    \le~& \frac{\kappa}{m} \tr \Rectbrac{ \Pi_\bot \Brac{ \int \nabla \nabla_j\matT \log \ell(\bx) \, \nabla_j \nabla\matT \log \ell(\bx)\, \pi'(\bx) \mdd \bx } \Pi_\bot\matT } \\
    =~&\frac{\kappa}{m} \tr \Rectbrac{ \Pi_\bot \widetilde{\sfR}_j \Pi_\bot\matT },
\end{split}
\end{equation}
where we denote $\widetilde{\sfR}_j := \int \nabla \nabla_j\matT \log \ell(\bx)\, \nabla_j \nabla\matT \log \ell(\bx)\, \pi'(\bx) \mdd \bx$. 
Here the second inequality follows from $\Pi_{j,r}\matT \Pi_{j,r}\preceq \kappa I_{d_j}$, since $\Pi_{j,r} = \sfC_{j}^{\frac{1}{2}} \widetilde{\Pi}_{j,r} \sfC_{j}^{-\frac{1}{2}} $ for some Euclidean orthogonal projection $\widetilde{\Pi}_{j,r} := \sfC_j^{-\frac{1}{2}} \sfU_j^{} \sfU_j^{\matT} \sfC_j^{-\frac{1}{2}} $, and thus
\begin{align*}
    \forall v \in \mR^{d_j},\quad v\matT \Pi_{j,r}\matT \Pi_{j,r} v =~& v\matT \sfC_{j}^{\frac{1}{2}} \widetilde{\Pi}_{j,r} \sfC_j^{-1} \widetilde{\Pi}_{j,r} \sfC_{j}^{-\frac{1}{2}} v \le \norm{\sfC_j^{-1}} \norm{ \widetilde{\Pi}_{j,r} \sfC_{j}^{\frac{1}{2}} v }^2 \\
    \le~& \norm{ \sfC_j^{-1}} \norm{ \sfC_{j}^{\frac{1}{2}} v }^2 \le \norm{C_j^{-1}} \norm{ \sfC_{j} } \norm{ v }^2 \le \kappa \norm{v}^2.
\end{align*}
Here we use the assumption that $m I \preceq \sfC \preceq M I$, and as a principle submatrix of $\sfC$, we have $m I \preceq \sfC_j \preceq M I$ and thus $\norm{C_j^{-1}} \norm{ C_{j} } \le \kappa$.


\ \\[-5pt]
\noindent
(3) For the last term, note that
\[
    \pi_0(\widetilde{\bx}_\bot|\bx_r) = \frac{1}{Z_0} \exp \Brac{ - \frac{1}{2} \normo{ \widetilde{\bx}_\bot - \sfC_{\bot,r} \sfC_{r,r}^\dag \bx_r }^2_{\sfC_{\bot|r}}  }.
\]
Here we denote 
\[
    \sfC_{\bot,r} = \Pi_{\bot} \sfC \Pi_{r}, \quad \sfC_{r,r} = \Pi_{r} \sfC \Pi_{r}, \quad \sfC_{\bot|r} = \sfC_{\bot,\bot} - \sfC_{\bot,r} \sfC_{r,r}^\dag \sfC_{r,\bot}.
\]
Since $\pi_0(\widetilde{\bx}_\bot|\bx_r)$ depends on $\widetilde{\bx}_\bot,\bx_r$ through the term $\widetilde{\bx}_\bot - \sfC_{\bot,r} \sfC_{r,r}^\dag \bx_r$, we have
\[
    \nabla_{\bx_r} \pi_0(\widetilde{\bx}_\bot|\bx_r) = - \sfC_{r,r}^\dag \sfC_{r,\bot} \nabla_{\bx_\bot} \pi_0(\widetilde{\bx}_\bot|\bx_r). 
\]
Denote $\sfK = \sfC_{r,r}^\dag \sfC_{r,\bot}$. By integration by parts, we have 
\begin{align*}
    \int \log \ell(\bx_r,\widetilde{\bx}_\bot) \nabla_{\bx_r} \pi_0(\widetilde{\bx}_\bot|\bx_r) \mdd \widetilde{\bx}_\bot =~& - \sfK \int \log \ell(\bx_r,\widetilde{\bx}_\bot) \nabla_{\bx_\bot} \pi_0(\widetilde{\bx}_\bot|\bx_r) \mdd \widetilde{\bx}_\bot \\
    =~& \sfK \int \nabla_{\bx_\bot} \log \ell(\bx_r,\widetilde{\bx}_\bot) \pi_0(\widetilde{\bx}_\bot|\bx_r) \mdd \widetilde{\bx}_\bot.
\end{align*}
%
The $j$-th block of the above inequality reads (note $\nabla_{\bx_{j,r}} f = \Pi_{j,r} \nabla_j f$ and $\nabla_{\bx_{k,\bot}} f = \Pi_{k,\bot} \nabla_k f$)
\begin{align*}
    & \norm{ \int \log \ell(\bx_r,\widetilde{\bx}_\bot) \Pi_{j,r} \nabla_j \pi_0(\widetilde{\bx}_\bot|\bx_r) \mdd \widetilde{\bx}_\bot } \\
    \le~& \sum_{k\in[K]} \norm{\sfK(j,k)} \norm{ \int \Pi_{k,\bot} \nabla_k \log \ell(\bx_r,\widetilde{\bx}_\bot) \pi_0(\widetilde{\bx}_\bot|\bx_r) \mdd \widetilde{\bx}_\bot } \\
    \le~& \sum_{k\in[K]} \norm{\sfK(j,k)} \int \norm{ \Pi_{k,\bot} \nabla_k \log \ell(\bx_r,\widetilde{\bx}_\bot) }\pi_0(\widetilde{\bx}_\bot|\bx_r) \mdd \widetilde{\bx}_\bot. 
\end{align*}
Since $\pi'(\bx) = \pi'(\bx_\bot|\bx_r) \pi'(\bx_r) = \pi_0(\bx_\bot|\bx_r) \pi'(\bx_r)$, we have
\begin{align*}
    \sfI_3 =~& \int \norm{ \int \log \ell(\bx_r,\widetilde{\bx}_\bot) \Pi_{j,r} \nabla_j \pi_0(\widetilde{\bx}_\bot|\bx_r) \mdd \widetilde{\bx}_\bot } \pi'(\bx) \mdd \bx \\
    \le~& \int \sum_{k\in[K]} \norm{\sfK(j,k)} \Brac{ \int \norm{ \Pi_{k,\bot} \nabla_k \log \ell(\bx_r,\widetilde{\bx}_\bot) } \pi_0(\widetilde{\bx}_\bot|\bx_r) \mdd \widetilde{\bx}_\bot } \pi'(\bx_r) \mdd \bx_r \\
    =~& \sum_{k\in[K]} \norm{\sfK(j,k)} \int \norm{ \Pi_{k,\bot} \nabla_k \log \ell(\bx) } \pi'(\bx) \mdd \bx. 
\end{align*}
%
As in the control of $\sfI_1$, this term can be controlled by
\[
    \sfI_3 \le \sum_{k\in[K]} \norm{\sfK(j,k)} \Brac{ \tr \Rectbrac{ \Pi_{k,\bot} \sfR_k \Pi_{k,\bot}\matT } }^{1/2}.
\]
By \Cref{lem:Kl1Bound}, $\sum_{k\in[K]} \norm{\sfK(j,k)} \le S(1+S) \nu! \kappa^{\nu+2}$. Therefore,
\begin{equation}
\label{eqn:pf_LLIS_I3}
    \sfI_3 \le S(1+S) \nu! \kappa^{\nu+2}\, \max_{k\in[K]} \Brac{ \tr \Rectbrac{ \Pi_{k,\bot} \sfR_k \Pi_{k,\bot}\matT } }^{1/2}.
\end{equation}
Substituting \eqref{eqn:pf_LLIS_I1}, \eqref{eqn:pf_LLIS_I2}, and \eqref{eqn:pf_LLIS_I3} into \eqref{eqn:pf_LLIS_err_decomp}, we have
\begin{align*}
    & \max_{j\in[K]} \norm{ \nabla_j \log \pi - \nabla_j \log \pi' }_{L^1(\pi')} \\
    \le~& \Brac{ 1 + S(1+S) \nu! \kappa^{\nu+2} } \max_{k\in[K]} \Brac{ \tr \Rectbrac{ \Pi_{k,\bot} \sfR_k \Pi_{k,\bot}\matT } }^{1/2} + \Brac{\frac{\kappa}{m}}^{1/2} \max_{k\in[K]} \Brac{ \tr \Rectbrac{ \Pi_\bot \widetilde{\sfR}_k \Pi_\bot\matT } }^{1/2}.
\end{align*}
Therefore, by the marginal transport inequality, we have 
\[
    \max_{i\in[K]} \sfW_1(\pi_i,\pi_i') \le C \Brac{ \max_{j\in[K]} \tr \Rectbrac{ \Pi_{j,\bot} \sfR_j \Pi_{j,\bot}\matT } + \max_{j\in[K]}  \tr \Rectbrac{ \Pi_\bot \widetilde{\sfR}_j \Pi_\bot\matT } }^{1/2},
\]
for some dimension-independent constant $C$ that depends on $S,\nu,\kappa,m$ only.
\end{proof}

\end{document}
